% !TEX program = lualatex
\documentclass[12pt, a4paper]{article}
\usepackage{master}

\DeclareWorkGR{Cat}{범주론}{Κατηγορίαι}{Categoriae}
\DeclareWorkGR{DI}{명제론}{Περὶ ἑρμηνείας}{De interpretatione}

\fancyhead[L]{\footnotesize\color{midgray}오르가논 강독}
\fancyhead[R]{\footnotesize\color{midgray}제 10 주}

\begin{document}
\thispagestyle{firstpage}

\weeksection{제10주}{명제론 17b16--19a22}{대당 사각형, 숨겨진 복합성, 해전 논변}

\subsection{대당 사각형의 완성 (7장 후반, 17b16--18a12)}

\subsubsection{네 가지 명제 형식과 네 가지 관계}

지난 주 7장 전반부에서 대당 사각형의 기초를 놓았습니다. 17b16--26에서 아리스토텔레스는 네 가지 명제 형식 사이의 대립 관계를 체계적으로 서술합니다. 중세 라틴 전통에서 \gri{affirmo}(긍정)의 모음 A, I와 \gri{nego}(부정)의 모음 E, O로 표기되었습니다. 네 가지 관계는: 모순(\gr{ἀντιφατικαί}, A/O 및 E/I) --- 정확히 하나가 참; 반대(\gr{ἐναντίαι}, A/E) --- 둘 다 참일 수 없으나 둘 다 거짓일 수 있음; 소반대(I/O) --- 둘 다 거짓일 수 없으나 둘 다 참일 수 있음; 함축(A→I, E→O).

\subsubsection{현대의 비판과 파슨스의 방어}

대당 사각형에 대한 현대의 비판(퍼스 1880, 우카시에비치 1929): S가 공집합일 때, I가 거짓이면 그 모순인 E가 참이어야 하고, E의 함축인 O도 참이어야 하는데, 이것은 존재 전제를 침해합니다.\footnote{SEP, ``The Traditional Square of Opposition'' (Parsons, rev.\ 2025).}

파슨스(Parsons)의 방어는 결정적입니다: 아리스토텔레스의 실제 O-형식은 `모든 S가 P인 것은 아니다'(\gr{οὐ πᾶς})이지, `어떤 S는 P가 아니다'가 아닙니다. S가 공집합일 때, A는 거짓이고(긍정 명제는 존재 전제를 가짐), O는 자명하게 참입니다. 비정합성은 보에티우스가 O-형식을 바꿔 표현한 데서 발생했으며, 아리스토텔레스의 원래 체계는 2천 년 이상 정합적이었습니다.

\subsection{숨겨진 복합성 (8장, 18a13--27)}

\subsubsection{명제는 언제 진정으로 단일한가?}

18a13--17에서 아리스토텔레스는 묻습니다: 명제가 진정으로 단일(\gr{μία})하려면 무엇이 필요한가? 문제는 동음이의(\gr{ὁμωνυμία})에서 발생합니다: `외투'(\gr{ἱμάτιον})가 `말'과 `사람' 모두를 의미한다면, `외투는 흰다'는 단일한 진술이 아니라 두 진술입니다(18a18--26). 밥지엔(Bobzien 2008)은 8장이 변증술적 논쟁에서의 동음이의 문제를 다루는 것이라고 주장합니다.\footnote{Susanne Bobzien, ``Aristotle's \gri{De Interpretatione} 8 is About Ambiguity,'' in \gri{Maieusis: Essays in Honour of Myles Burnyeat}, ed.\ D.\ Scott (Oxford: Oxford University Press, 2008).}

\subsubsection{RCP의 두 번째 예외}

동음이의적 명제에서는 겉보기 모순 쌍이 둘 다 거짓일 수 있습니다(18a26--27). 휘태커(1996)는 이것을 RCP의 두 번째 예외로 해석합니다: 7장의 한정 없는 명제에 이어, 8장의 동음이의적 명제도 규칙에서 벗어납니다. 8장은 9장으로의 다리 역할을 합니다: RCP가 보편적 모순 쌍과 단칭 명제에 대해 성립하지만, 한정 없는 명제(7장)와 동음이의적 명제(8장)에서는 예외가 있음을 확인한 뒤, 이제 제3의 예외 --- 미래 우연 단칭 명제(9장) --- 가 있는지를 묻습니다.

\subsection{해전 논변 (9장, 18a28--19a22)}

\subsubsection{문제의 설정과 결정론적 논변}

18a28--33에서 아리스토텔레스는 요약합니다: 과거·현재 진술과 보편적으로 양화된 보편자에 대해서는 RCP가 성립합니다. `그러나 미래의 개별자들(\gr{τῶν καθ' ἕκαστα καὶ μελλόντων})에 대해서는 사정이 다릅니다.'

18a34--18b25에서 결정론적 논변이 전개됩니다. `내일 해전이 있을 것이다'가 지금 참이라면, 해전은 필연적입니다 --- 참으로 진술된 것은 일어나지 않을 수 없습니다. 일반화(18b5--9): `일어나는 모든 것은 필연적으로 일어난다' --- 전면적 결정론. 숙고와 행위는 무의미합니다(18b26--19a6). 아리스토텔레스는 이 결론을 거부합니다: `우리는 미래에 올 것이 숙고와 행위에 기원을 가진다는 것을 본다'(18b31). 외투의 예(19a12--16): `이 외투는 잘릴 수도 있고 잘리지 않을 수도 있다; 그러나 실제로는 잘리지 않고 먼저 닳아서 없어질 것이다' --- 이것은 진정한 우연성을 보여줍니다.

\subsubsection{두 가지 학술적 해석}

명제론 9장에 대한 학술적 해석은 크게 두 가지로 나뉩니다.\footnote{SEP, ``Medieval Theories of Future Contingents'' (Knuuttila, rev.\ 2024).}

(A) 전통적 해석: 아리스토텔레스는 미래 우연 명제에 대해 이가 원리를 제한합니다 --- 사건이 발생하기 전까지 이 명제들은 (확정적으로) 참도 거짓도 아닙니다. 그러나 배중률은 보존됩니다. 주요 지지자: 우카시에비치(1920/1957), 프레데(1985), 소라브지(1980), 크리벨리(2004), 바이데만(2012). 프레데(1985)는 가장 영향력 있는 현대적 옹호를 제공했습니다: 아리스토텔레스는 미래 우연 모순 쌍에서 진리와 거짓이 `확정적으로'(\gr{ἀφωρισμένως}) 분배되지 않는다고 부정합니다.\footnote{Dorothea Frede, ``The Sea-Battle Reconsidered,'' \gri{Oxford Studies in Ancient Philosophy} 3 (1985): 31--87.}

(B) 비표준적 해석: 아리스토텔레스는 이가 원리를 부정하지 않습니다. 미래 우연 명제는 참이거나 거짓이지만, 그 참이 확정적이거나 필연적인 참이 아닙니다. 주요 지지자: 앤스콤(1956), 힌티카(1973), 파인(1984), 휘태커(1996). 앤스콤은 아리스토텔레스의 관심사가 참의 유무가 아니라 참의 양태라고 주장합니다.\footnote{G.\ E.\ M.\ Anscombe, ``Aristotle and the Sea Battle,'' \gri{Mind} 65 (1956): 1--15.}

\subsubsection{사고 연습 --- 해전 논변과 범주론의 반대자 수용}

3주차에서 다룬 `반대자 수용'(\gr{δεκτικὸν τῶν ἐναντίων}) --- 실체의 가장 고유한 특성 --- 을 떠올려 봅시다. 소크라테스가 지금 앉아 있고 나중에 서 있을 수 있다면, `소크라테스는 앉아 있다'는 지금 참이고 나중에 거짓이 됩니다. 명제의 진리값이 시간에 따라 변한다는 이 아이디어는 해전 논변의 배경에 있습니다. 범주론 5장에서 이미 `진술이 참에서 거짓으로 바뀔 때 변하는 것은 무엇인가?'라는 물음이 제기되었습니다 --- 명제론 9장은 바로 그 물음의 가장 급진적인 전개입니다.

\subsection*{참고문헌}
\begin{references}

\refitem Anscombe, G.\ E.\ M. ``Aristotle and the Sea Battle.'' \gri{Mind} 65 (1956): 1--15.

\refitem Bobzien, Susanne. ``Aristotle's \gri{De Interpretatione} 8 is About Ambiguity.'' In \gri{Maieusis}, edited by D.\ Scott. Oxford: Oxford University Press, 2008.

\refitem Crivelli, Paolo. \gri{Aristotle on Truth}. Cambridge: Cambridge University Press, 2004.

\refitem Frede, Dorothea. ``The Sea-Battle Reconsidered.'' \gri{Oxford Studies in Ancient Philosophy} 3 (1985): 31--87.

\refitem Sorabji, Richard. \gri{Necessity, Cause, and Blame}. London: Duckworth, 1980.

\refitem Whitaker, C.\ W.\ A. \gri{Aristotle's De Interpretatione: Contradiction and Dialectic}. Oxford: Clarendon Press, 1996.

\end{references}

\end{document}
